\documentclass{article}

\usepackage{amsmath,amssymb}
\usepackage{xurl}
\usepackage[hidelinks]{hyperref}
\setlength{\emergencystretch}{2em}

% Shared commands for notes in docs/paper-gaps/.

\DeclareUnicodeCharacter{2080}{\ifmmode{}_{0}\else\textsubscript{0}\fi}
\DeclareUnicodeCharacter{2081}{\ifmmode{}_{1}\else\textsubscript{1}\fi}
\DeclareUnicodeCharacter{2082}{\ifmmode{}_{2}\else\textsubscript{2}\fi}
\DeclareUnicodeCharacter{2099}{\ifmmode{}_{n}\else\textsubscript{n}\fi}

\newcommand{\Lean}{\textsc{Lean}}
\ExplSyntaxOn
\NewDocumentCommand{\leanid}{m}{
  \ifmmode
    \textsf{\detokenize{#1}}
  \else
    \group_begin:
    \urlstyle{sf}
    \tl_set:Nn \l_tmpa_tl {#1}
    \tl_replace_all:Nnn \l_tmpa_tl {\_} {_}
    \exp_args:NV \path \l_tmpa_tl
    \group_end:
  \fi
}
\ExplSyntaxOff
\providecommand{\tr}{\operatorname{tr}}
\newcommand{\tnleanIssueBase}{https://github.com/LionSR/TNLean/issues/}
\newcommand{\tnleanPrBase}{https://github.com/LionSR/TNLean/pull/}
\newcommand{\tnleanDocBase}{https://sirui-lu.com/TNLean/docs/}
\newcommand{\ghissue}[1]{\href{\tnleanIssueBase#1}{GitHub issue \##1}}
\newcommand{\ghpr}[1]{\href{\tnleanPrBase#1}{PR \##1}}
\newcommand{\doclink}[1]{\href{\tnleanDocBase#1.html}{\path{#1}}}

% Dirac notation and the identity operator, as in blueprint/src/macros/common.tex.
% \providecommand keeps note-local definitions authoritative.
\usepackage{dsfont}
\providecommand{\ket}[1]{|#1\rangle}
\providecommand{\bra}[1]{\langle #1|}
\providecommand{\braket}[2]{\langle #1 | #2 \rangle}
\providecommand{\ketbra}[2]{|#1\rangle\!\langle #2|}
\providecommand{\Id}{\mathds{1}}
\providecommand{\one}{\mathds{1}}
\providecommand{\C}{\mathbb{C}}
\providecommand{\MN}[1]{M_{#1}(\C)}
\providecommand{\MD}{M_D(\C)}

% External referencing for paper sources.  The build helper writes sanitized
% auxiliary files under build/paper-aux/ and excludes duplicated source labels.
\usepackage{xr}

% Import a generated paper auxiliary file with a caller-supplied label prefix.
% The candidate paths support builds from docs/paper-gaps/, the repository root,
% and build/paper-aux/ itself.
\newcommand{\importpaperaux}[2]{%
  \IfFileExists{../../build/paper-aux/#2.aux}{%
    \externaldocument[#1]{../../build/paper-aux/#2}%
  }{%
    \IfFileExists{build/paper-aux/#2.aux}{%
      \externaldocument[#1]{build/paper-aux/#2}%
    }{%
      \IfFileExists{#2.aux}{%
        \externaldocument[#1]{#2}%
      }{}%
    }%
  }%
}

% General external reference macros
\newcommand{\paperref}[2]{\ref{#1-#2}}
\newcommand{\papereqref}[2]{(\ref{#1-#2})}

% CPSV16 source (1606.00608 / MPDO-22-12-17-2.tex) external document and shortcuts
\importpaperaux{cpsv16-}{1606.00608/MPDO-22-12-17-2-tnlean}
\newcommand{\cpsvref}[1]{\paperref{cpsv16}{#1}}
\newcommand{\cpsveqref}[1]{\papereqref{cpsv16}{#1}}

% Verdict marker: \gapnote{<kind>}{<status>}. Kind is the policy
% classification (clarification, local-correction, scope-restriction,
% unfaithful, false-source, open-gap); status is open, wip (elimination
% underway), resolved, or historical. Machine-read by the site index;
% prints a verdict line.
\newcommand{\gapnote}[2]{\par\noindent\textbf{Verdict:} #1 (#2).\par}


\title{The parent commuting Hamiltonian condition in Cirac--Perez-Garcia--Schuch--Verstraete 2016}
\author{}
\date{2026-06-15}

\begin{document}
\maketitle
\gapnote{scope-restriction}{open}

\section{Source condition}

The condition in Appendix~D of Ref.~\cite{Cirac2016MPDO_arXiv} is not the route
used to prove the forward RFP-to-NNCPH implication. The proof of
Theorem~3.10 at source lines~1305--1307 consists of one sentence stating that
the implication follows immediately from the structural characterization in
Theorem~3.11 of Ref.~\cite{Cirac2016MPDO_arXiv}. Definition~D.2 belongs to the
later, independent discussion of decorrelation.

Definition~D.2 of Ref.~\cite{Cirac2016MPDO_arXiv} considers a tripartite
Hilbert space $\mathcal H_A\otimes\mathcal H_X\otimes\mathcal H_B$ and local
projectors $Q_{AX}$ and $Q_{XB}$. Let $K_{AX}$ and $K_{XB}$ denote their
kernels. Their lifts to the tripartite space are
\begin{align}
    \widehat Q_{AX}
        &=Q_{AX}\otimes I_B,
    \label{eq:cpsv16_parent_commuting_hamiltonian_scope_qax_lift}\\
    \widehat Q_{XB}
        &=I_A\otimes Q_{XB}.
    \label{eq:cpsv16_parent_commuting_hamiltonian_scope_qxb_lift}
\end{align}
The source condition is\footnote{Local source:
\path{Papers/1606.00608/MPDO-22-12-17-2.tex}:2205--2218, labels
\path{QAXQXB} and \path{KAXBKAXKXB}.}
\begin{align}
    [\widehat Q_{AX},\widehat Q_{XB}]
        &=0,
    \label{eq:cpsv16_parent_commuting_hamiltonian_scope_source_commutation}\\
    K_{AXB}
        &=(K_{AX}\otimes\mathcal H_B)
          \cap(\mathcal H_A\otimes K_{XB}).
    \label{eq:cpsv16_parent_commuting_hamiltonian_scope_source_kernel_intersection}
\end{align}

The proof of Proposition~D.3 of Ref.~\cite{Cirac2016MPDO_arXiv} introduces
complementary projectors and their lifts:
\begin{align}
    P_{AX}
        &=I_{AX}-Q_{AX},
    \label{eq:cpsv16_parent_commuting_hamiltonian_scope_pax_complement}\\
    P_{XB}
        &=I_{XB}-Q_{XB},
    \label{eq:cpsv16_parent_commuting_hamiltonian_scope_pxb_complement}\\
    \widehat P_{AX}
        &=P_{AX}\otimes I_B,
    \label{eq:cpsv16_parent_commuting_hamiltonian_scope_pax_lift}\\
    \widehat P_{XB}
        &=I_A\otimes P_{XB}.
    \label{eq:cpsv16_parent_commuting_hamiltonian_scope_pxb_lift}
\end{align}
These lifted projectors satisfy
\begin{align}
    \widehat P_{AX}\widehat P_{XB}
        &=P_{AXB}.
    \label{eq:cpsv16_parent_commuting_hamiltonian_scope_projector_product}
\end{align}
This is the product identity for commuting orthogonal projectors onto the two
local kernels.\footnote{Local source:
\path{Papers/1606.00608/MPDO-22-12-17-2.tex}:2259--2289, in particular labels
\path{PAXQAX}, \path{PXBQXB}, and \path{arrggg}.}

The explicit tripartite declaration
\leanid{TripartiteDecorrelation.CommutingParentHamiltonian} records the
range-intersection equality. For commuting orthogonal projectors,
\begin{align}
    \operatorname{ran}P_{AXB}
        =\operatorname{ran}\widehat P_{AX}
          \cap\operatorname{ran}\widehat P_{XB}
        \quad&\Longleftrightarrow\quad
        \widehat P_{AX}\widehat P_{XB}=P_{AXB}.
    \label{eq:cpsv16_parent_commuting_hamiltonian_scope_range_product_equivalence}
\end{align}
This equivalence is proved by
\leanid{TripartiteDecorrelation.hasGroundSpaceIntersection_iff_product_eq};
its forward implication is exposed as
\leanid{TripartiteDecorrelation.CommutingParentHamiltonian.hproduct}.

\section{Current boundary}

The declaration \leanid{Decorrelation.HasCommutingParentHam} records only the
idempotent-product part of the source condition.\footnote{See
\doclink{TNLean/MPS/ParentHamiltonian/Decorrelation}.} For a complex vector
space $E$, it asks for endomorphisms $P_{AX},P_{XB}:E\to E$ satisfying
\begin{align}
    P_{AX}^{2}
        &=P_{AX},
    \label{eq:cpsv16_parent_commuting_hamiltonian_scope_abstract_pax_idempotent}\\
    P_{XB}^{2}
        &=P_{XB},
    \label{eq:cpsv16_parent_commuting_hamiltonian_scope_abstract_pxb_idempotent}\\
    P_{AX}P_{XB}
        &=P_{XB}P_{AX},
    \label{eq:cpsv16_parent_commuting_hamiltonian_scope_abstract_commutation}\\
    P_{AX}P_{XB}
        &=P_K.
    \label{eq:cpsv16_parent_commuting_hamiltonian_scope_abstract_product}
\end{align}
It does not assert a tensor decomposition of $E$, local action on the $AX$ and
$XB$ factors, self-adjointness, orthogonality, or the existence of source
projectors $Q_{AX}$ and $Q_{XB}$. Every idempotent $P_K$ therefore admits the
tautological choice
\begin{align}
    P_{AX}
        &=P_K,
    \label{eq:cpsv16_parent_commuting_hamiltonian_scope_tautological_pax}\\
    P_{XB}
        &=P_K,
    \label{eq:cpsv16_parent_commuting_hamiltonian_scope_tautological_pxb}
\end{align}
recorded by \leanid{Decorrelation.HasCommutingParentHam.ofIdem}. Thus this
declaration is not Definition~D.2 of Ref.~\cite{Cirac2016MPDO_arXiv}; it is
the algebraic consequence sufficient for the proved backward decorrelation
implication.

A separate three-site support layer records the $AX$ and $XB$ faces:
\leanid{MPSTensor.leftPairLift}, \leanid{MPSTensor.rightPairLift}, and
\leanid{MPSTensor.HasOverlappingTwoSiteCommutation}. The source
kernel-intersection condition is represented by
\begin{align}
    K_{AXB}
        &=\ker((Q_{AX})_{AX})\cap\ker((Q_{XB})_{XB})
    \label{eq:cpsv16_parent_commuting_hamiltonian_scope_lifted_kernel_intersection}
\end{align}
in \leanid{MPSTensor.HasAppendixD2ParentCommutingHamiltonian}. Its forgetful
projection gives the algebraic overlapping-commutation predicate.

The scalar pair-index relation has been assembled into the physical map
\begin{align}
    (Uv)_i
        &=\sum_{\alpha,\beta}U^i_{\alpha,\beta}v_{\alpha,\beta},
    \label{eq:cpsv16_parent_commuting_hamiltonian_scope_physical_map}\\
    U^*U
        &=\Id.
    \label{eq:cpsv16_parent_commuting_hamiltonian_scope_physical_isometry}
\end{align}
These identities are recorded in
\leanid{MPSTensor.AppendixBStructuralData.physicalIsometry} and
\leanid{MPSTensor.AppendixBStructuralData.physicalIsometryLeftInverse_comp}.
For every $L$, the tensor-power maps satisfy
\begin{align}
    U^{*\otimes L}U^{\otimes L}
        &=\Id,
    \label{eq:cpsv16_parent_commuting_hamiltonian_scope_tensor_power_isometry}
\end{align}
as recorded by
\leanid{MPSTensor.AppendixBStructuralData.physicalIsometryTensorPowerLeftInverse_comp}.

On two sites, the inserted virtual bond is
\begin{align}
    (I_\varphi v)_{(a,b),(b',c)}
        &=\delta_{b,b'}\lambda_bv_{a,c}.
    \label{eq:cpsv16_parent_commuting_hamiltonian_scope_virtual_bond_embedding}
\end{align}
Its physical image satisfies
\begin{align}
    \operatorname{ran}(U^{\otimes 2}I_\varphi)
        &=G_2(A_{\mathrm{core}}).
    \label{eq:cpsv16_parent_commuting_hamiltonian_scope_two_site_image}
\end{align}
This is
\leanid{MPSTensor.AppendixBStructuralData.twoSiteBasicEmbedding_range}. The
orthogonal projector onto this image is
\leanid{MPSTensor.AppendixBStructuralData.twoSiteBasicSupportProjection}; its
complement equals $q_2(A_{\mathrm{core}})$ by
\leanid{MPSTensor.AppendixBStructuralData.twoSiteBasicSupportProjection_eq_complement}.

The passage to complementary projectors is recorded by
\leanid{MPSTensor.HasOverlappingTwoSiteCommutation.complement}. Three-site
transport is available directly for $Q_{AX},Q_{XB}$ through
\leanid{MPSTensor.localTerm_two_three_zero_one_commute_of_appendixD2}, and for
their complements through
\leanid{MPSTensor.localTerm_two_three_zero_one_commute_of_appendixD2_complement}.

For the Appendix~B structural datum, the coefficient-space representatives are
\leanid{MPSTensor.AppendixBStructuralData.appendixBQAXOnCoeffSpace} and
\leanid{MPSTensor.AppendixBStructuralData.appendixBQXBOnCoeffSpace}. Both are
identified with $q_2(A)$ before placement on the $AX$ and $XB$ faces. Under
the canonical reindexings, their factor actions are
\begin{align}
    [\widehat Q_{AX}]_{AX}
        &=[\widehat Q_{AX}]\otimes I_d,
    \label{eq:cpsv16_parent_commuting_hamiltonian_scope_qax_factor_action}\\
    [\widehat Q_{XB}]_{XB}
        &=I_d\otimes[\widehat Q_{XB}].
    \label{eq:cpsv16_parent_commuting_hamiltonian_scope_qxb_factor_action}
\end{align}
The generic identities are
\leanid{MPSTensor.leftPairLift_toMatrix_reindex} and
\leanid{MPSTensor.rightPairLift_toMatrix_reindex}; their Appendix~B
specializations are
\leanid{MPSTensor.AppendixBStructuralData.appendixBQAX_factor_action} and
\leanid{MPSTensor.AppendixBStructuralData.appendixBQXB_factor_action}.

This construction realizes only the homogeneous choice
\begin{align}
    \mathcal H_A
        &=\mathbb C^d,
    \label{eq:cpsv16_parent_commuting_hamiltonian_scope_homogeneous_a}\\
    \mathcal H_X
        &=\mathbb C^d,
    \label{eq:cpsv16_parent_commuting_hamiltonian_scope_homogeneous_x}\\
    \mathcal H_B
        &=\mathbb C^d,
    \label{eq:cpsv16_parent_commuting_hamiltonian_scope_homogeneous_b}
\end{align}
introduced at source lines~2185--2186. It does not provide the arbitrary
factorization used in Definition~D.2 at source lines~2205--2218 of
Ref.~\cite{Cirac2016MPDO_arXiv}.

Because $q_2(A)$ is the canonical orthogonal projector onto $G_2(A)^\perp$,
the coefficient-space maps are orthogonal. They do not by themselves
construct the source projectors on
$\mathcal H_A\otimes\mathcal H_X$ and
$\mathcal H_X\otimes\mathcal H_B$, as stated in
\path{TNLean/MPS/RFP/AppendixBStructuralData.lean}. The required transported
identification is an explicit hypothesis of
\leanid{MPSTensor.AppendixBStructuralData.hasOverlappingTwoSiteCommutation_of_appendixD2_identified_projectors}.

Without this identification, commutation follows from the adjacent
virtual-bond projectors and their isometric transport, as recorded by
\leanid{MPSTensor.AppendixBStructuralData.twoSiteBasicSupportProjection_commute_lifts}
and
\leanid{MPSTensor.AppendixBStructuralData.hasOverlappingTwoSiteCommutation}.
The adjacent parent terms commute on every periodic chain with $N>2$ by
\leanid{MPSTensor.AppendixBStructuralData.adjacent_twoSite_localTerms_commute}.

For an injective tensor, the three-site intersection identity is
\begin{align}
    G_3(A)
        &=\ker(q_2(A)_{AX})\cap\ker(q_2(A)_{XB}).
    \label{eq:cpsv16_parent_commuting_hamiltonian_scope_three_site_intersection}
\end{align}
It is recorded by
\leanid{MPSTensor.groundSpace_three_eq_adjacent_twoSite_parent_kernels}.
Consequently, the Appendix~B representatives satisfy the coefficient-space
analogues of the idempotence, commutation, and kernel-intersection conditions,
with
\begin{align}
    \widehat K_{AXB}
        &=G_3(A).
    \label{eq:cpsv16_parent_commuting_hamiltonian_scope_coefficient_ground_space}
\end{align}
This is proved by
\leanid{MPSTensor.AppendixBStructuralData.hasAppendixD2ParentCommutingHamiltonian}.
The theorem does not identify the hatted maps with the source projectors on
the explicitly factored Hilbert spaces. For a normal left-canonical RFP tensor,
the RFP hypotheses supply the required injectivity, as proved by
\leanid{MPSTensor.rfp_hasAppendixD2ParentCommutingHamiltonian_of_leftCanonical}.

\section{Elimination plan}

The general finite-dimensional tripartite equivalence is recorded by
\leanid{TripartiteDecorrelation.parentHamiltonian_iff_decorrelated}. Its parent
condition uses orthogonal projectors on the explicit factors and states their
lifted range intersection. The theorem
\leanid{TripartiteDecorrelation.hasGroundSpaceIntersection_iff_product_eq}
equates this condition with
(\ref{eq:cpsv16_parent_commuting_hamiltonian_scope_projector_product}). The
forward direction constructs the projectors as supports of the two marginals;
the reverse direction derives and uses the product identity. The local
correction to the marginal-support expansion at source lines~2236--2252 is
documented in
\path{docs/paper-gaps/cpsv16_decorrelation_slice_expansion_fix.tex}.

The coefficient-space analogues are proved. The remaining task is to construct
the source projectors on the stated tensor factors and identify their
transported actions with the coefficient-space maps. The conditional theorem
\leanid{MPSTensor.AppendixBStructuralData.hasOverlappingTwoSiteCommutation_of_appendixD2_identified_projectors}
states precisely this missing identification.

This comparison belongs to the RFP application, not to the general
decorrelation equivalence. The declaration
\leanid{Decorrelation.HasCommutingParentHam} remains only the algebraic
idempotent-product consequence, and
\leanid{MPSTensor.HasAppendixD2ParentCommutingHamiltonian} does not itself
record self-adjointness. The explicit tripartite condition does record
self-adjointness and proves both directions of Proposition~D.3 of
Ref.~\cite{Cirac2016MPDO_arXiv}.

The remaining obligations for Theorem~3.10 of
Ref.~\cite{Cirac2016MPDO_arXiv} are the full canonical-form scope and the
all-chain ground-space spanning statement. They are tracked by
\ghissue{2633}. Definition~D.2 should not be presented as their missing proof
route.

\bibliographystyle{alpha}
\bibliography{references}

\end{document}
